\section{Введение}
\label{sec:Chapter0} \index{Chapter0}

\subsection{Введение в область}
\label{subsec:intro-area}

Большие языковые модели (Large Language Models, LLM) за последние годы совершили качественный скачок в решении задач обработки естественного языка: генерации текста, ответов на вопросы, суммаризации и машинного перевода~\cite{brown2020language, touvron2023llama, openai2023gpt4}. Одним из ключевых направлений развития современных LLM стала способность к \textit{многошаговым рассуждениям} (multi-hop reasoning) --- последовательному выводу ответа через цепочку логически связанных промежуточных шагов. Техники, основанные на цепочке рассуждений (Chain-of-Thought, CoT)~\cite{wei2022chain}, продемонстрировали, что побуждение модели к пошаговому размышлению существенно повышает точность на задачах, требующих логического и арифметического вывода.

Тем не менее, несмотря на впечатляющие результаты, цепочки рассуждений LLM по-прежнему подвержены характерным ошибкам: галлюцинациям фактов, нарушению логической связности между шагами и накоплению ошибок при увеличении глубины вывода. Эти проблемы особенно выражены в задачах, требующих интеграции разрозненных фактов из обширной базы знаний, --- так называемых \textit{multi-hop} вопросах. Стандартные подходы к оценке качества рассуждений, опирающиеся на бинарное сравнение финального ответа с эталонным, не позволяют диагностировать \textit{структурные} дефекты в промежуточных шагах вывода.

Графы знаний (Knowledge Graphs, KG) представляют собой структурированное хранилище фактов в виде троек $(\text{субъект}, \text{отношение}, \text{объект})$ и давно применяются для задач извлечения информации, вопросно-ответных систем и верификации фактов. Крупномасштабные графы знаний, такие как Wikidata, содержат миллионы сущностей и отношений, образуя богатую фактическую основу для оценки корректности рассуждений. Использование графовой структуры для \textit{явного} представления эталонных цепочек вывода позволяет перейти от бинарной оценки <<правильно/неправильно>> к \textit{дифференцированной} оценке качества каждого шага рассуждения.

\subsection{Актуальность}
\label{subsec:relevance}

Актуальность исследования обусловлена несколькими факторами. Во-первых, растущая потребность в надёжных и верифицируемых рассуждениях LLM для применения в критически важных областях (медицина, право, научные исследования) требует перехода от оценки результата к оценке процесса. Во-вторых, графовые структуры предоставляют естественный формализм для представления логических зависимостей, однако их систематическое использование для обучения и оценки LLM остаётся слабо изученным. В-третьих, компактные модели (порядка 0.6--7B параметров) представляют особый практический интерес ввиду доступности для дообучения и развёртывания, но их способность к структурно корректным рассуждениям существенно уступает более крупным моделям.

\subsection{Исследование существующих решений}
\label{subsec:existing-solutions}

Существующие подходы к улучшению многошаговых рассуждений LLM можно разделить на несколько категорий.

\textbf{Промптинговые методы.} Семейство методов Chain-of-Thought (CoT)~\cite{wei2022chain} и его расширения --- Self-Consistency~\cite{wang2023selfconsistency}, Tree-of-Thought~\cite{yao2023tree}, Graph-of-Thought~\cite{besta2024graph} --- направлены на улучшение рассуждений на этапе инференса путём структурирования процесса генерации. Эти методы не требуют дообучения модели, однако не устраняют фундаментальные ограничения, связанные с качеством внутренних представлений модели.

\textbf{Retrieval-Augmented Generation (RAG).} Подходы RAG~\cite{lewis2020retrieval} дополняют генерацию LLM релевантными документами из внешнего корпуса, что снижает частоту галлюцинаций. Графовые расширения RAG (Graph-RAG) используют графы знаний для более точного извлечения контекста. Однако данные методы решают проблему доступа к фактам, но не обучают модель корректному логическому выводу.

\textbf{Обучение с подкреплением для рассуждений.} Методы RLHF~\cite{ouyang2022training} и их развитие (GRPO~\cite{shao2024deepseekmath}, RLOO) позволяют обучать модели на основе функций награды, оценивающих качество генерации. Большинство существующих подходов используют бинарную награду (правильный/неправильный финальный ответ) или процессную награду (Process Reward Models, PRM)~\cite{lightman2023lets}, обучаемую на человеческих аннотациях отдельных шагов. Ключевой недостаток PRM --- высокая стоимость получения обучающих данных и субъективность аннотаций.

\textbf{Интеграция графов знаний в LLM.} Ряд работ исследует использование графов знаний для повышения фактической точности LLM~\cite{pan2024unifying}: внедрение графовых эмбеддингов в архитектуру модели, использование KG для верификации фактов, генерация структурированных объяснений. Тем не менее, использование графовой структуры непосредственно в качестве \textit{сигнала обучения} при RL-дообучении остаётся практически неизученным направлением.

Таким образом, существующие решения либо ограничиваются этапом инференса (промптинг, RAG), либо требуют дорогостоящих человеческих аннотаций (PRM), либо не используют структурную информацию графов знаний для обучения моделей корректным рассуждениям. Данная работа предлагает заполнить этот пробел.

\subsection{Цель работы}
\label{subsec:goal}

Целью данной работы является исследование влияния явной графовой структуры на способность больших языковых моделей выполнять многошаговые рассуждения, а также разработка методов оценки и обучения, учитывающих структуру reasoning-процесса.

Основная идея подхода заключается в следующем: для заданного вопроса в графе знаний выделяется эталонный подграф, описывающий корректную цепочку логического вывода; ответ модели преобразуется в аналогичное графовое представление; затем вычисляются метрики структурного соответствия между сгенерированным и эталонным графами. Полученная мера соответствия используется в качестве функции награды при дообучении модели методами обучения с подкреплением, что позволяет не только оценивать правильность финального ответа, но и поощрять корректность промежуточных шагов рассуждения.

Для достижения поставленной цели решаются следующие \textbf{задачи}:
\begin{enumerate}
    \item Разработка программного комплекса для автоматизированного преобразования текстовых генераций LLM в графовое представление и сопоставления с эталонными графами рассуждений.
    \item Формализация метрик структурного соответствия графов (precision, recall, F1, Graph Edit Distance) для количественной оценки качества цепочек рассуждений.
    \item Сравнительный анализ рассуждающих способностей ряда моделей (Qwen3, Llama~3.x, OLMo-Think, GPT-OSS) на датасетах RuleTaker и Wikidata.
    \item Реализация двухэтапного пайплайна обучения (SFT~+~RL) с использованием графовой функции награды, учитывающей структурное соответствие сгенерированного подграфа эталонному.
    \item Экспериментальная проверка гипотезы о том, что интеграция графовой награды повышает точность логических выводов и снижает долю структурных разрывов в цепочках рассуждений.
\end{enumerate}

% \textbf{Структура работы.} В~Главе~\ref{sec:Chapter1} приводится обзор существующих подходов к многошаговым рассуждениям LLM и использованию графов знаний. В~Главе~\ref{sec:Chapter2} описывается постановка задачи и предлагаемая методология. В~Главе~\ref{sec:Chapter3} представлена архитектура разработанного программного комплекса. В~Главе~\ref{sec:Chapter4} приводятся результаты экспериментов. В~Главе~\ref{sec:Chapter5} формулируются выводы и направления дальнейших исследований.